\section{Materials and Methods}
\label{sec:methods}

\subsection{Dataset: Corpus, Annotation and Splits}
\label{sec:dataset}

The corpus is 32 nocturnal road-transect videos recorded with a vehicle-mounted FLIR thermal
camera at four sites in southern Illinois, eight videos per site. Frames are $640\times512$
single-channel thermal at 60\,fps. The vehicle travels at 25--90\,km/h, so the background
translates continuously and the same animal is visible for a few seconds at most.
\Cref{fig:dataset} shows representative frames with their annotation and \cref{tab:dataset}
summarizes the corpus. The top row is what a detector receives: in the third and fourth panels the animals are a few dozen pixels of slightly warmer texture against thermal clutter. The median animal is $29\times24$\,px, roughly the smallest box in
the third panel. Thermal transects carry low and drifting global contrast, so every frame is
normalized with contrast-limited adaptive histogram equalisation
(CLAHE)~\cite{reza2004realization} before use, identically at training and inference and for
every model reported here. This single preprocessing choice is worth more than any
architectural one we tested: test mAP50 falls from $0.519$ to $0.299$ without it. The figures
in this paper therefore show normalized rather than raw frames, because they are the pixels
the detectors consume.

\begin{figure}[htbp]
\centering
\setlength{\tabcolsep}{1.5pt}
\renewcommand{\arraystretch}{0.6}
\begin{tabular}{cccc}
\includegraphics[width=0.228\textwidth]{figures/dataset/p1_raw.jpg} &
\includegraphics[width=0.228\textwidth]{figures/dataset/p2_raw.jpg} &
\includegraphics[width=0.228\textwidth]{figures/dataset/p3_raw.jpg} &
\includegraphics[width=0.228\textwidth]{figures/dataset/p4_raw.jpg} \\
\includegraphics[width=0.228\textwidth]{figures/dataset/p1_gt.jpg} &
\includegraphics[width=0.228\textwidth]{figures/dataset/p2_gt.jpg} &
\includegraphics[width=0.228\textwidth]{figures/dataset/p3_gt.jpg} &
\includegraphics[width=0.228\textwidth]{figures/dataset/p4_gt.jpg} \\[2pt]
\small (\textbf{a}) two ranges & \small (\textbf{b}) close approach &
\small (\textbf{c}) group, three ranges & \small (\textbf{d}) mid-range pair \\
\end{tabular}
\caption{\textbf{Dataset examples.} Top row: raw thermal frames after the CLAHE
normalization the detector is trained and evaluated with. Bottom row: the same frames with
their ground-truth boxes. Panels show the \emph{large} end of the corpus for legibility at
print scale.\label{fig:dataset}}
\end{figure}

\begin{table}[htbp]
\caption{\textbf{Dataset composition.} 32 thermal road-transect videos from four Illinois
sites, $640\times512$ at 60\,fps, a single annotated class. Splits are by video and
site-stratified. Rows count the frames \emph{extracted} for detector training, not the
521,930 frames of the corpus; the complete track annotation is 21,646 boxes, of which the
12,903 here fall inside the extracted frames.\label{tab:dataset}}
\centering
\begin{tabularx}{\textwidth}{lRRRRR}
\toprule
\textbf{Split} & \textbf{Videos} & \textbf{Frames} & \textbf{With Deer} & \textbf{Empty} & \textbf{Boxes} \\
\midrule
Train & 19 & 12,852 & 5013 & 7839 & 7547 \\
Val   &  4 &  2337 &   935 & 1402 & 2564 \\
Test  &  9 &  3725 & 1490 & 2235 & 2792 \\
\midrule
\textbf{Total} & \textbf{32} & \textbf{18,914} & \textbf{7438} & \textbf{11,476} & \textbf{12,903} \\
\bottomrule
\end{tabularx}
\end{table}

Two frame counts appear in this paper and they measure different things. The 32 videos
contain 521,930 frames, and \emph{every} deer in \emph{every} frame is annotated in the
Computer Vision Annotation Tool (CVAT) as a member of an individual track. That annotation is
the 236 tracks and 21,646 boxes quoted above, and it is what ``fully annotated'' means here.
At 60\,fps consecutive frames are near-duplicates, so the detector dataset of \cref{tab:dataset} is an extraction from the annotation: 18,914 frames sampled around annotated activity together
with explicit negatives, carrying the 12,903 boxes that fall inside them. The 21,646 figure
is the size of the ground truth and the 12,903 the size of the detector's training signal.
The counting evaluation of \cref{sec:counting} always runs against the former, on whole
videos decoded end to end. Of the 18,914 extracted frames, 7438 contain at least one deer
and 11,476 are explicit negatives; background frames matter, since a survey system must not
fire on warm rocks or structures.

The 236 animals are unevenly distributed across the four sites: SHB contributes 132 (56\% of
the corpus), TON 51, SHW 38 and MAS 15. One TON track is marked outside the frame on every
frame, the animal being occluded throughout, so it can never be detected. \textbf{235 animals
are therefore detectable, and every detection and counting denominator in this paper is
235}, while 236 remains the count of annotated individuals. The imbalance is not incidental.
SHB is the dense-group site where most of the counting error concentrates
(\cref{sec:counting}), and MAS is small enough that its per-site figures should be read as
indicative only. Every individual deer is annotated as a distinct \emph{track}, and this is
the property the paper depends on. One track is one animal, so the same annotation serves as
per-frame detection ground truth, as association ground truth and as count ground truth,
the number of tracks in a video. In \cref{sec:reid} it also serves as identity labels.

We split by \emph{video}, never by frame, and stratify by site so that all four sites appear
in each split. The deployment target is these same sites, and a frame-level split would leak
near-duplicate frames across the boundary. The result is 19 training videos (12,852
images), 4 validation (2337) and 9 test (3725), allocated by animal count to roughly
65/19/16\%. Throughout the paper, \textbf{``held out'' means the 13 validation and test
videos the detector never trained on}, containing 83 of the 236 deer. \Cref{sec:heldout}
shows why the distinction matters: counting accuracy measured over all 32 videos, 19 of
which trained the detector, is 15.5 points higher than on held-out video alone.

\subsection{Object Scale}
\label{sec:scale}

The animals are small and the distribution is one-sided. The median box is $29\times24$
pixels ($\sqrt{\text{area}}=26.8$), 71\% of boxes fall below the COCO~\cite{lin2014coco}
small-object threshold of $32^2$, and \emph{zero} boxes reach the COCO-large category.
Resolution is therefore the strongest single accuracy lever, which motivates the ablation in
\cref{sec:detection}. The one-sidedness has a second consequence, developed in
\cref{sec:analysis}: no box in the training split exceeds 96 pixels, and detector
performance falls on the few animals that appear larger. The upper tail of the scale
distribution is as under-served as the lower one, and it is far less often checked.

\subsection{Counting Pipeline}
\label{sec:pipeline}

\begin{figure}[htbp]
  \centering
  \includegraphics[width=\textwidth]{figures/wildlife_pipeline.png}
  \caption{\textbf{The TRACT counting pipeline.} Thermal frames are contrast-normalized,
detected per frame, associated into candidate tracks, and filtered by a three-parameter
  confirmation rule; the count is the number of accepted tracks. Dashed components were
  ablated and are not used. Frames are illustrative.\label{fig:pipeline}}
\end{figure}

\Cref{fig:pipeline} shows the system, which we call \textbf{TRACT} (Thermal Road-transect
Animal Counting and Tracking). Every component is standard; the assembly of contrast normalisation, detection, motion-compensated tracking, orphan recovery and a calibrated confirmation rule is what is arranged for counting rather than for detection. Detections from YOLO11m at 640\,px are
associated by BoT-SORT~\cite{aharon2022botsort} with sparse-optical-flow global motion
compensation, which is mandatory here because the whole frame translates as the vehicle
moves. The result is a set of candidate tracks. One point of terminology belongs here,
because the association stage is easily mistaken for something it is not. \textbf{The
reported system runs with BoT-SORT's appearance term disabled and does not attempt to
recognise individual animals.} It associates detections into tracks within a single video
using motion and geometry, which prevents an animal from being counted twice in one
transect. It could not recognise individuals in any case, since \cref{sec:reid} shows that
at a median $29\times24$ pixels a two-number size descriptor out-performs a pretrained
appearance backbone. We call this \emph{track association} throughout. Where appearance
appears in this paper it is an \emph{ablation} (pool G of \cref{sec:inversion}) or a probe
of how much identity signal the imagery carries (\cref{tab:reid}), never a component of the
reported counts.

Only the detector is gradient-trained. The tracker's Kalman filter, motion compensation and
association carry no weights, and the confirmation rule's three parameters are fitted by
grid search. The pipeline also has no stochastic component at inference. Run three times on
the same GPU over four held-out videos holding 55 of the 83 held-out animals, it produced
byte-identical detections, tracks and counts. A count is therefore a fixed function of the
video, and repeating a run cannot reduce its error. Across platforms, one benchmarked
detector evaluated on an NVIDIA GH200, with a different processor architecture and different
PyTorch and CUDA builds, reproduced its A100 keyframe AP$_{50}$ of 0.670; counts themselves
were compared only on identical hardware. The detector is not the one with the best
AP$_{50}$, and the choice is deliberate. \Cref{sec:choosing} selects on the
precision-weighted $F_{0.5}$ instead. The confirmation stage must reject tens of thousands
of false candidates against roughly two hundred true ones, so it pays far more for a false
positive than for a false negative. That section also verifies the choice end to end: the
highest-AP$_{50}$ model run through this same pipeline counts fewer animals.

The confirmation stage decides which candidate tracks are real, distinct animals, and the
count is the number it accepts. Let a candidate track $\tau$ consist of $m(\tau)$
detections spanning $s(\tau)$ seconds, with per-frame detector confidences ordered as
$c_{(1)}\ge c_{(2)}\ge\cdots\ge c_{(m)}$. Its track score is the mean of its $k$ highest
confidences,
\begin{equation}
\bar{c}_k(\tau)=\frac{1}{k}\sum_{j=1}^{k} c_{(j)},\qquad k=5,
\label{eq:topk}
\end{equation}
and the rule accepts $\tau$ when three conditions hold at once,
\begin{equation}
a(\tau)=\ind\!\left[m(\tau)\ge m^{*}\right]\,\ind\!\left[s(\tau)\ge s^{*}\right]\,
\ind\!\left[\bar{c}_5(\tau)\ge c^{*}\right],
\qquad
\hat{c}_v=\sum_{\tau\in\mathcal{T}_v} a(\tau),
\label{eq:rule}
\end{equation}
where $\ind[\cdot]$ is the indicator function, $\mathcal{T}_v$ is the candidate set of video
$v$ and $\hat{c}_v$ is its count. The score is a property of the \emph{track}: a
low-confidence frame is never discarded, it continues to extend the track, and a candidate
is rejected only when its five best frames average below $c^{*}$. The three thresholds are
not set by hand. They minimise counting error over the 19 detector-training videos
$\mathcal{V}_{\mathrm{fit}}$,
\begin{equation}
(m^{*},s^{*},c^{*})=\operatorname*{arg\,min}_{(m,s,c)\in\mathcal{G}}\;
\frac{1}{|\mathcal{V}_{\mathrm{fit}}|}\sum_{v\in\mathcal{V}_{\mathrm{fit}}}
\bigl|\hat{c}_v(m,s,c)-N_v\bigr|,
\label{eq:gridsearch}
\end{equation}
where $N_v$ is the annotated number of animals in $v$ and the grid $\mathcal{G}$ is
$m\in\{1,2,3,5,8,12,20\}$ frames, $s\in\{0,0.1,0.25,0.5,1,2\}$\,s and
$c\in\{0,0.15,0.25,\dots,0.85\}$, 378 cells in all. The values are frozen before any
held-out video is scored (\cref{sec:heldout}). The search selects $m^{*}=20$, $s^{*}=0$ and
$c^{*}=0.65$. At that operating point the span condition imposes nothing, since $s\ge0$
accepts a track of any duration, so the rule is in effect two parameters. We keep $s$
because it is the parameter the search uses to compensate when $c$ is lowered
(\cref{sec:sensitivity}), and \cref{fig:rulegrid} shows the whole search. We additionally
recover \emph{orphan} detections, those the tracker returns without a track identity, by
linking detections within a temporal gap and a scale-normalized distance into pseudo-tracks.
Without this, single-frame animals are discarded entirely; with it, 17 animals are
recovered.

\subsection{Calibrated Per-Animal Confidence}
\label{sec:calibration}

We calibrate a per-track
posterior for ``this track is on a real animal'' and report it with every counted animal.
\textbf{92.1\% of counted animals carry confidence $\ge0.80$} (mean 0.965, area under the
ROC curve 0.893), and 793 low-confidence tracks are flagged for human review rather than
auto-counted. The calibrator is \textbf{isotonic regression fitted on the validation videos
only}~\cite{zadrozny2002transforming}: a monotone step function $p=\phi(\bar{c}_5)$ that
maps the confirmer's raw score to a probability by least squares under a monotonicity
constraint. No held-out video contributes to the fit. We prefer isotonic to Platt scaling
because the raw score is not sigmoidal in the log-odds, and a monotone step function makes
no shape assumption. The cost is that the output takes a finite number of levels, 46 of
them on held-out video. We also propagate the per-track probabilities into an uncertainty
on each video's count. If the accepted tracks of video $v$ carry probabilities
$p_1,\dots,p_{\hat{c}_v}$ and are treated as independent, the number of real animals among
them is Poisson-binomial~\cite{hong2013poisson}, with
\begin{equation}
\mathrm{Var}\bigl[\hat{c}_v\bigr]=\sum_{i=1}^{\hat{c}_v} p_i\,(1-p_i).
\label{eq:pb}
\end{equation}
That quantity is narrower than it looks. It expresses uncertainty about whether the tracks
already accepted are real animals and says nothing about animals that were never detected
or never accepted. On held-out video its nominal 95\% band therefore contains the true
count in only one of thirteen videos. It should not be read as an interval on abundance,
which would need a model of the probability that an animal is detected and counted at all.

\Cref{tab:calib} reports whether the calibrated number behaves like a probability on videos
the calibrator never saw. A track scores as positive when it is the \emph{primary} track of a
real animal; duplicates count as negatives, which is the decision the counting stage makes.
We use the two standard summaries~\cite{guo2017calibration}. With the $n$ scored tracks
partitioned into $B=10$ equal-count bins $\mathcal{B}_b$ by predicted probability, the
expected calibration error~\cite{naeini2015obtaining} and the Brier
score~\cite{brier1950verification} are
\begin{equation}
\mathrm{ECE}=\sum_{b=1}^{B}\frac{|\mathcal{B}_b|}{n}\,\bigl|\,\overline{y}_b-\overline{p}_b\bigr|,
\qquad
\mathrm{BS}=\frac{1}{n}\sum_{i=1}^{n}\bigl(p_i-y_i\bigr)^2,
\label{eq:ece}
\end{equation}
where $y_i\in\{0,1\}$ is the label, $p_i$ the calibrated probability, and $\overline{y}_b$
and $\overline{p}_b$ the mean label and mean probability in bin $b$. Over all 2211 held-out
candidate tracks the ECE is 0.004 and the Brier score 0.017, against 0.029 for a constant
predictor at the base rate. That aggregate is generous, because 74\% of candidates sit at
the isotonic floor where no decision is ever taken. Restricted to the 316 tracks in the
decision region the ECE is 0.034 and the Brier score 0.093 against 0.148. The second is the one a practitioner should read. One caveat belongs with them:
isotonic regression saturates, so the top level reports 1.000 on 13 held-out tracks whose
true positive rate is 0.846. Confidence at the very top of the scale is optimistic by
roughly 15 points, and the flag-for-review band, not the top of the scale, is where the
calibrated number earns its keep.

\begin{table}[htbp]
\caption{\textbf{Reliability of the per-track posterior} on the 13 held-out videos, against
the isotonic calibrator fitted on validation videos only. Positives are primary tracks;
duplicates are negatives. The constant-predictor column is the Brier score of always
predicting the base rate, so the gap is the skill the posterior
adds.\label{tab:calib}}
\centering\footnotesize
\begin{tabularx}{\textwidth}{LCCCCC}
\toprule
\textbf{Scope} & \textbf{Tracks} & \textbf{Base rate} & \textbf{ECE} & \textbf{Brier} & \textbf{Constant} \\
\midrule
All held-out candidates & 2211 & 0.030 & 0.004 & 0.017 & 0.029 \\
Decision region ($p>0.01$) & 316 & 0.180 & 0.034 & 0.093 & 0.148 \\
\midrule
Fit videos, for reference & 4797 & 0.029 & 0.003 & 0.017 & 0.029 \\
\bottomrule
\end{tabularx}
\end{table}

Raw detector confidence cannot serve this purpose. No track in the corpus exceeds 0.90 and
the median counted track sits at 0.74, because the animals are 27 pixels, so the requirement
is only meaningful against a calibrated posterior. \Cref{fig:evidence} shows why a score
must be per track rather than per frame. The raw confidence varies by 0.15 across frames of
the same animal (Animal B, 0.66--0.81) even where every frame is a clean detection, so no
single frame's value characterises the track. Every counted animal is also exported as a contact sheet of its own track, one panel per sampled frame with elapsed time and per-frame score, so a human can verify any count against the frames that produced it.

\begin{figure}[htbp]
  \centering
  \includegraphics[width=0.86\linewidth]{figures/evidence_two.jpg}
  \caption{\textbf{Per-animal evidence, two individuals.} Each row is one animal's own
  track, one panel per sampled frame, with elapsed time and raw per-frame detector confidence
  beneath. Both animals are from videos the detector never trained on.\label{fig:evidence}}
\end{figure}

\subsection{Matching Criterion}
\label{sec:criterion}

A detection metric asks whether a predicted box \emph{localizes} an animal; a survey asks
whether the animal was \emph{noticed}. Intersection over union (IoU), the standard
localization score~\cite{everingham2010pascal,lin2014coco}, diverges sharply from the survey
question on a 27-pixel target, and the divergence can be made exact.

\begin{proposition}[IoU under a centre offset]
\label{prop:iou}
Let a predicted box and a ground-truth box be axis-aligned squares of equal side $s$ whose
centres differ by $d<s$ pixels along one axis. Then
\begin{equation}
\mathrm{IoU}(d)=\frac{s-d}{s+d},
\qquad\text{so that}\qquad
\mathrm{IoU}(d)\ge t \iff d\le s\,\frac{1-t}{1+t}.
\label{eq:iou}
\end{equation}
\end{proposition}
\begin{proof}
The intersection is a rectangle of sides $s$ and $s-d$, with area $s(s-d)$. The union is
the sum of the two areas minus the intersection, $2s^{2}-s(s-d)=s(s+d)$. The ratio gives
the first equality, and solving $(s-d)/(s+d)\ge t$ for $d$ gives the second.
\end{proof}

At the corpus median $s=27$, \cref{eq:iou} says that a centre error of more than 9 pixels
fails IoU$\ge$0.50 and one of more than 3.9 pixels fails IoU$\ge$0.75. A box offset by six
pixels is therefore scored as both a false positive and a false negative at IoU 0.5, while
for counting it has correctly found the animal. We therefore evaluate under four matching
rules: IoU$\geq$0.50, IoU$\geq$0.30, centre-in-box and \emph{any overlap}. The last is the
counting criterion, reported as presence or counting recall and never as mAP, with the
strict criterion alongside it everywhere. The criterion is validated against
centre-in-box, an independent counting-oriented rule that accepts a prediction only when
its centre falls inside the animal and so rejects oversized boxes outright. The two agree
to within 0.032 F1 across all eleven models, so the conclusions rest on two criteria
agreeing rather than on one being chosen. Precision is reported alongside recall
throughout, which further constrains box quality.

\subsection{Decomposing Counting Error}
\label{sec:decomposition}

Let $G_v$ be the ground-truth animals in video $v$, $N_v=|G_v|$, and $\mathcal{T}_v$ the
candidate tracks the pipeline produces. For each animal $g\in G_v$ we define three
indicators:
\begin{itemize}[leftmargin=*,labelsep=4.5mm]
\item $R(g)=1$ (\reachedq) if $g$ is touched, on at least one frame, by at least one
candidate. If an animal is not reached, no downstream stage can recover it; this measures
\textbf{detection and tracking}.
\item $P(g)=1$ (\primaryq) if, among the candidates touching $g$, the one that covers $g$ on
the most frames is not also the best cover for some other animal. Two animals sharing a
single best-covering candidate are indistinguishable to any rule that emits one count per
track; this measures \textbf{association}.
\item $K(g)=1$ (\countedq) if the confirmation stage accepted a candidate covering $g$.
This measures \textbf{confirmation}, and it is the end-to-end result.
\end{itemize}
The stage rates are the means of these indicators over the animals under evaluation,
$\rho=\overline{R}$, $\pi=\overline{P}$ and $\kappa=\overline{K}$. An accepted candidate
is a candidate, so $K\le R$ and $P\le R$ by construction, and in every configuration we
report $\kappa\le\pi\le\rho$, so the three read as a funnel whose steps attribute error to
the three stages. We report all three for every configuration; reporting only the last, as
is conventional, makes a detection miss and a confirmation miss look identical.

The video-level error has an exact decomposition of its own. Credit each accepted track to
at most one animal, the one it covers on the most frames, and each animal to at most one
accepted track, and let $A_v$ be the number of animals credited in video $v$. A track that
covers two animals credits only one of them, so the pair contributes one credited animal and
one miss. Write $D_v=\hat{c}_v-A_v\ge0$ for the accepted tracks that credit no new animal
(duplicates of an animal already credited, and tracks on no animal) and $M_v=N_v-A_v\ge0$ for
the animals with no credited track.

\begin{proposition}[Signed error is duplicates minus misses]
\label{prop:identity}
For every video, $\hat{c}_v-N_v=D_v-M_v$, and
$\min(\hat{c}_v,N_v)=A_v+\min(D_v,M_v)$.
\end{proposition}
\begin{proof}
Both follow from $\hat{c}_v=A_v+D_v$ and $N_v=A_v+M_v$. Subtracting gives the first
identity; taking the minimum of the two and factoring out $A_v$ gives the second.
\end{proof}

\Cref{prop:identity} is the bookkeeping the rest of the paper relies on. A count is unbiased
only when the duplicates it admits equal the animals it misses, which is the trade
\cref{sec:sensitivity} measures, and every intervention that recovers animals is judged by
whether it lowers $M_v$ faster than it raises $D_v$. The aggregate metrics over a set of
videos $\mathcal{V}$ are
\begin{equation}
\mathrm{MAE}=\frac{1}{|\mathcal{V}|}\sum_{v\in\mathcal{V}}\bigl|\hat{c}_v-N_v\bigr|,
\qquad
\mathrm{Bias}=\frac{1}{|\mathcal{V}|}\sum_{v\in\mathcal{V}}\bigl(\hat{c}_v-N_v\bigr),
\qquad
\kappa_{\mathrm{cap}}=\frac{\sum_{v}\min(\hat{c}_v,N_v)}{\sum_{v}N_v},
\label{eq:metrics}
\end{equation}
where $\kappa_{\mathrm{cap}}$ is the \emph{counted} fraction, capped per video so that
over-counting in one video cannot mask misses in another. By \cref{prop:identity} it exceeds
the identity-matched fraction $\sum_v A_v/\sum_v N_v$ by exactly $\sum_v\min(D_v,M_v)/\sum_v
N_v$. The capped aggregate thus credits a video for predicting the right \emph{number} even
when some accepted tracks do not land on the animals, and \cref{sec:counting} shows the two
differ by 10 points. Counting quantities carry 95\% intervals throughout. Proportions over
the 83 held-out animals use the Wilson score interval~\cite{wilson1927probable}, and means
over the 13 held-out videos use a $t$-interval on the per-video values. These are sampling
intervals over the corpus, and with 83 animals they are wide: a three-animal difference is
not resolvable. We deliberately do \emph{not} attach intervals to the detection metrics of
\cref{sec:detection}. Each architecture was trained once, so a bootstrap over the test
images would give a narrow interval that describes test-set sampling while saying nothing
about the training variance that separates one architecture from another.

\subsection{Held-Out Protocol}
\label{sec:heldout}

Counting runs process all 32 videos, but the detector trained on 19 of them, and the
confirmation rule is conventionally swept on the same videos it is reported on. Both leaks
inflate the result. Our protocol closes both: \textbf{sweep the rule on the 19
detector-training videos, freeze it, and report on the 13 held-out videos}. Learned
confirmers are trained and thresholded under the identical split, so the comparison in
\cref{sec:learned} is like-for-like.

\begin{table}[htbp]
\caption{\textbf{Why the held-out protocol matters.} Evaluating on all 32 videos, as an in-sample pipeline evaluation would, reports 81.8\% coverage. On video the detector never trained on
it is 66.3\%. Both figures use the capped per-video fraction of \cref{sec:decomposition}; the
uncapped ratio of predicted to true animals reads 89.0\% and 69.9\% respectively, and is the
more flattering of the two because over-counting one video offsets misses in another.\label{tab:leak}}
\centering
\begin{tabularx}{\textwidth}{lRRR}
\toprule
\textbf{Evaluation Scope} & \textbf{Videos} & \textbf{Deer} & \textbf{Coverage} \\
\midrule
Detector-training videos & 19 & 153 & 90.2\% \\
\textbf{Held out (val + test)} & \textbf{13} & \textbf{83} & \textbf{66.3\%} \\
All 32 (the default) & 32 & 236 & 81.8\% \\
\bottomrule
\end{tabularx}
\end{table}

The protocol is worth 15.5 points (\cref{tab:leak}): counting accuracy measured over all 32
videos is 81.8\%, and on the 13 held-out videos it is 66.3\%. Four of the 13 held-out videos
are the detector's \emph{validation} split. No gradient was computed on them, but they were
used for checkpoint selection, so they are held out in the sense that matters for the
confirmation rule and only partly held out for the detector. They carry 45 of the 83
animals. We state this rather than reporting on the 9 test videos alone, for two reasons.
The confirmation rule, the subject of the paper, saw none of the 13, and splitting the
held-out set further would leave 38 animals to carry every counting conclusion.

\subsection{Box Size as a Range Proxy}
\label{sec:rangeproxy}

The corpus records no range to the animals, so box size stands in for it in
\cref{sec:analysis,sec:operating}. The relation follows from the pinhole
camera~\cite{hartley2004multiple}: an animal of physical length $L$ at range $r$ subtends
\begin{equation}
s=\frac{fL}{r}
\label{eq:pinhole}
\end{equation}
pixels, where $f$ is the focal length in pixels. Box size is therefore inversely
proportional to range for animals of similar size, and it conflates range with body size and
posture, a caveat that applies wherever the proxy is used. The proxy also yields a null
expectation for how many animals should be annotated at each size, which
\cref{sec:operating} uses to test the ground truth itself for range-dependent censoring.

\begin{proposition}[Availability law for box size]
\label{prop:availability}
Suppose animals of length $L$ are distributed with uniform density $\lambda$ per unit ground
area inside the camera's horizontal field of view of angle $\theta$, and every animal in view
is annotated. Then the expected number of annotated animals with box size in $[a,b)$ is
\begin{equation}
E[a,b)=\frac{\lambda\theta\,(fL)^{2}}{2}\left(\frac{1}{a^{2}}-\frac{1}{b^{2}}\right),
\label{eq:availability}
\end{equation}
so the density of animals per unit box size falls as $s^{-3}$. If instead the animals lie
in a strip of fixed width $w$ along the road, $E[a,b)=\lambda w fL\,(1/a-1/b)$, a density
falling as $s^{-2}$.
\end{proposition}
\begin{proof}
The ground in view between ranges $r$ and $r+dr$ is a circular sector of area
$\theta r\,dr$, holding $\lambda\theta r\,dr$ animals in expectation. By \cref{eq:pinhole},
$r=fL/s$ and $|dr|=fL\,s^{-2}\,ds$, so the expected number with box size in $[s,s+ds)$ is
$\lambda\theta (fL)^{2}s^{-3}\,ds$. Integrating from $a$ to $b$ gives
\cref{eq:availability}. In a strip of width $w$ the ground between $r$ and $r+dr$ has area
$w\,dr$, and the same substitution gives $\lambda w fL\,s^{-2}\,ds$, whose integral is
$\lambda w fL\,(1/a-1/b)$.
\end{proof}

The constant in front is unknown, but it cancels in ratios. Fitting it to the number of
animals observed in a reference size range predicts the number expected in every other
range, which is how \cref{sec:operating} applies the result. Distance sampling makes the
same uniform-density assumption for animals relative to a transect
line~\cite{buckland2015distance,thomas2010distance}, and a road transect can violate
it~\cite{diefenbach2025accounting}, which is why the strip law is carried along as the more
conservative alternative.
